\documentclass[12pt,a4paper]{article}
\usepackage{multirow}
\usepackage[utf8]{inputenc}
\usepackage{geometry}
\geometry{margin=0.5in}
\usepackage{mathptmx} % Times New Roman font
\usepackage{helvet}
\usepackage{courier}
\usepackage{amsmath}
\usepackage{amsfonts}
\usepackage{amssymb}
\usepackage{graphicx}
\usepackage{booktabs}
\usepackage{array}
\usepackage{longtable}
\usepackage{url}
\usepackage{xcolor}
\usepackage{subcaption}
\usepackage{float}
\usepackage{algorithm}
\usepackage{algorithmic}
\usepackage{listings}
\usepackage{tcolorbox}
\tcbuselibrary{most}
\usepackage{hyperref}

% Professional color system
\definecolor{primaryblue}{RGB}{31,81,138}
\definecolor{secondaryblue}{RGB}{70,130,180}
\definecolor{accentgreen}{RGB}{40,167,69}
\definecolor{warningred}{RGB}{220,53,69}
\definecolor{highlightgray}{RGB}{248,249,250}
\definecolor{tableheader}{RGB}{233,236,239}
\definecolor{legacycolor}{RGB}{180,100,50}
\definecolor{neuralcolor}{RGB}{100,150,50}

% Professional highlighting commands
\newcommand{\primarytext}[1]{\textcolor{primaryblue}{\textbf{#1}}}
\newcommand{\secondarytext}[1]{\textcolor{secondaryblue}{\textbf{#1}}}
\newcommand{\success}[1]{\textcolor{accentgreen}{\textbf{#1}}}
\newcommand{\warning}[1]{\textcolor{warningred}{\textbf{#1}}}
\newcommand{\highlight}[1]{\colorbox{highlightgray}{#1}}
\newcommand{\legacy}[1]{\textcolor{legacycolor}{\textbf{#1}}}
\newcommand{\neural}[1]{\textcolor{neuralcolor}{\textbf{#1}}}

\lstset{
    basicstyle=\ttfamily\scriptsize,
    breaklines=true,
    frame=single,
    language=Python,
    backgroundcolor=\color{highlightgray}
}

\title{\textbf{Comprehensive Evaluation of Informed Contextual Multi-Armed Bandit Models for Quantum Network Routing}\\\
\large{Empirical Assessment of iCMAB Algorithms Across Stochastic and Adversarial Quantum Environments with Multi-Run and Multi-Capacity Analysis}}

\author{Graduate Research Report\\
Department of Computer Science\\
Rochester Institute of Technology\\
AI/Quantum Computing Research Track}

\date{December 11, 2025}

\begin{document}

\maketitle

\begin{abstract}
This report presents a comprehensive empirical evaluation of five informed contextual multi-armed bandit (iCMAB) algorithms specifically designed for quantum network routing under stochastic and adversarial conditions. Using multi-run experiment suites (3-run and 5-run protocols) across scaled capacity regimes (1$\times$, 1.5$\times$, 2$\times$) and evaluator type variations, we systematically evaluate: iCEpsilonGreedy, iCPursuit, iCThompsonSampling, iCEXP4, and iCEpochGreedy. Experimental scenarios span stochastic failures (Attack Rate: 0.25) and baseline optimal conditions, with performance measured as percentage efficiency relative to an oracle route. Results demonstrate that \primarytext{iCEpsilonGreedy decisively dominates all tested pure iCMAB scenarios}, achieving 86.8\% efficiency under stochastic conditions and 93.2\% under optimal baseline conditions. iCPursuit forms a consistent second tier at 68.5\% stochastic efficiency, with iCThompsonSampling providing a nearby baseline at 66.3\%, while iCEXP4 and iCEpochGreedy remain suboptimal (37.1\% and 37.6\%, respectively). For neural baselines, iCPursuitNeuralUCB achieves 88.9\% stochastic efficiency, GNeuralUCB achieves 85.2\%, and EXPNeuralUCB achieves 87.0\%, demonstrating multiple neural integration pathways for informed contextual methods. Across all environments and capacity scales, algorithms demonstrate consistent performance with variance $\leq$6.5\%, and statistical validation confirms 3-run suites are representative of 5-run behavior. These empirically grounded rankings provide critical baselines for future iCMAB-EXPNeuralUCB hybrid architectures and establish performance thresholds for hybrid development in quantum routing under realistic failure conditions.
\end{abstract}

\newpage
\tableofcontents

\newpage

% Executive Summary
\begin{tcolorbox}[colback=highlightgray,colframe=primaryblue,title=\textbf{Executive Summary}]
\begin{itemize}
\item \neural{Top Neural Performer}: iCPursuitNeuralUCB achieves 88.9\% Oracle efficiency (stochastic) with superior baseline performance of 91.1\%, exceptional stability (CV = 1.0\%), demonstrating strength of neural integration of informed pursuit methods.
\item \primarytext{Top Informed Performer}: iCEpsilonGreedy achieves 86.8\% Oracle efficiency (stochastic) with exceptional baseline performance of 93.2\% and solid variance (CV = 2.3\%) among top-performing pure iCMAB algorithms.
\item \neural{Secondary Neural}: GNeuralUCB achieves 85.2\% stochastic efficiency with exceptional consistency (CV = 0.8\%), providing Gaussian-style neural alternative.
\item \secondarytext{Performance Tiers}: A clear hierarchy emerges: Tier 1 -- iCPursuitNeuralUCB (88.9\% neural), iCEpsilonGreedy (86.8\% pure); Tier 1-alt -- GNeuralUCB (85.2\% neural), EXPNeuralUCB (87.0\% neural); Tier 2 -- iCPursuit (68.5\%) and iCThompsonSampling (66.3\%); Tier 3 -- iCEXP4 (37.1\%) and iCEpochGreedy (37.6\%).
\item \success{Robustness}: Coefficient of variation (CV) ranges from 0.6\% (iCEpochGreedy) to 6.5\% (iCPursuit) under stochastic conditions, with top performers maintaining CV $\leq$2.3\%.
\item \success{Capacity and Regime Stability}: Across scaled capacities (1$\times$--2$\times$) and different evaluator types, iCEpsilonGreedy and iCPursuitNeuralUCB maintain consistent ranking with variance $<$1\%, confirming dominance is fundamental rather than configuration-dependent.
\item \neural{Neural Consistency}: GNeuralUCB achieves 0.8\% CV stochastic consistency, superior to pure iCMAB methods except iCEpsilonGreedy, establishing neural reliability.
\item \success{Multi-Run Consistency}: 3-run and 5-run experiment suites yield differences $\leq$0.50\%, confirming that 3-run protocols already capture representative algorithm behavior.
\item \primarytext{Hybrid Integration Ready}: iCEpsilonGreedy (Tier 1 pure) and iCPursuitNeuralUCB (Tier 1 neural) are identified as primary candidates for iCMAB-EXPNeuralUCB hybrids, with explicit performance targets of $\geq$86.8\% for pure and $\geq$88.9\% for neural variants.
\item \warning{Critical Finding}: Prior hybrid development used iCPursuit (CMAB winner) without testing iCEpsilonGreedy (iCMAB winner) or neural variants. Switching to iCPursuitNeuralUCB offers an immediate 20.4\% performance gain (68.5\% $\rightarrow$ 88.9\%).
\end{itemize}
\end{tcolorbox}

\newpage
\section{Results and Analysis}

\begin{tcolorbox}[colback=tableheader,colframe=primaryblue,title=\textbf{Primary iCMAB Performance Results (Stochastic Environment, Attack Rate: 0.25)}]

\subsection{Global Performance Ranking}

\begin{table}[H]
\footnotesize
\centering
\caption{Global iCMAB and Neural Performance: Stochastic Failures (3-run, Multi-Capacity Aggregate)}
\label{tab:global-performance-stochastic}
\begin{tabular}{@{}lccccc@{}}\toprule
\textbf{Algorithm} & \textbf{Run 1 (\%)} & \textbf{Run 2 (\%)} & \textbf{Run 3 (\%)} & \textbf{Avg (\%)} & \textbf{CV (\%)} \\\midrule
\neural{iCPursuitNeuralUCB}   & \neural{87.8} & \neural{89.4} & \neural{89.5} & \neural{88.9} & \neural{1.0} \\
\primarytext{iCEpsilonGreedy}   & \primarytext{84.0} & \primarytext{88.0} & \primarytext{88.4} & \primarytext{86.8} & \primarytext{2.3} \\
\legacy{EXPNeuralUCB}           & \legacy{86.6} & \legacy{87.4} & \legacy{87.0} & \legacy{87.0} & \legacy{0.4} \\
\neural{GNeuralUCB}            & \neural{84.8} & \neural{85.4} & \neural{85.4} & \neural{85.2} & \neural{0.8} \\
iCPursuit                       & 70.6 & 62.3 & 72.6 & 68.5 & 6.5 \\
iCThompsonSampling              & 64.0 & 66.1 & 68.7 & 66.3 & 2.9 \\
iCEXP4                          & 37.4 & 37.1 & 36.7 & 37.1 & 0.8 \\
iCEpochGreedy                   & 37.8 & 37.7 & 37.3 & 37.6 & 0.6 \\\bottomrule
\end{tabular}
\end{table}

\textbf{Key Observations:}
\begin{itemize}
\item \neural{iCPursuitNeuralUCB Leadership}: Achieves 88.9\% average efficiency, leading all algorithms with exceptional consistency (CV = 1.0\%).
\item \neural{Neural Enhancement}: iCPursuitNeuralUCB outperforms its pure iCPursuit counterpart by 20.4 percentage points (68.5\% $\rightarrow$ 88.9\%), demonstrating substantial neural integration value.
\item \primarytext{iCEpsilonGreedy Strong Second}: Achieves 86.8\% average efficiency, outperforming second-place pure iCPursuit by 18.3 percentage points.
\item \primarytext{Tier 1 Excellence}: iCEpsilonGreedy's CV of 2.3\% indicates strong consistency across runs among pure algorithms.
\item \legacy{EXPNeuralUCB Competitive}: Achieves 87.0\% stochastic efficiency, positioned at tier-1 level with iCEpsilonGreedy and superior consistency (CV = 0.4\%).
\item \neural{GNeuralUCB Alternative}: Achieves 85.2\% stochastic efficiency with exceptional consistency (CV = 0.8\%), establishing Gaussian-style neural as competitive backup.
\item \secondarytext{Four-tier hierarchy}: (1) iCPursuitNeuralUCB (88.9\% neural); (2) iCEpsilonGreedy (86.8\% pure) + EXPNeuralUCB (87.0\% neural); (3) GNeuralUCB (85.2\% neural); (4) iCPursuit + iCThompsonSampling (68.5\%, 66.3\%); (5) iCEXP4 + iCEpochGreedy (37.1\%, 37.6\%).
\item \warning{iCPursuit Variability}: Despite pure-method second-place ranking, iCPursuit shows highest variance (CV = 6.5\%), indicating less stable performance.
\end{itemize}

\end{tcolorbox}

\subsection{Baseline Performance Analysis (Optimal Conditions, Attack Rate: 0.0)}

\begin{table}[H]
\centering
\caption{Baseline Performance: Optimal Conditions (No Attacks)}
\label{tab:baseline-performance}
\begin{tabular}{@{}lccccc@{}}\toprule
\textbf{Algorithm} & \textbf{Run 1 (\%)} & \textbf{Run 2 (\%)} & \textbf{Run 3 (\%)} & \textbf{Avg (\%)} & \textbf{CV (\%)} \\\midrule
\neural{iCPursuitNeuralUCB}   & \neural{90.8} & \neural{91.3} & \neural{91.3} & \neural{91.1} & \neural{0.3} \\
\primarytext{iCEpsilonGreedy}   & \primarytext{92.7} & \primarytext{93.4} & \primarytext{93.4} & \primarytext{93.2} & \primarytext{0.4} \\
\legacy{EXPNeuralUCB}           & \legacy{86.2} & \legacy{86.4} & \legacy{86.2} & \legacy{86.3} & \legacy{0.1} \\
\neural{GNeuralUCB}            & \neural{84.1} & \neural{84.7} & \neural{84.5} & \neural{84.4} & \neural{0.5} \\
iCPursuit                       & 66.2 & 70.4 & 71.1 & 69.2 & 3.5 \\
iCThompsonSampling              & 67.2 & 70.4 & 72.2 & 69.9 & 3.6 \\
iCEXP4                          & 38.8 & 38.0 & 39.3 & 38.7 & 1.9 \\
iCEpochGreedy                   & 39.6 & 39.2 & 39.0 & 39.3 & 0.5 \\\bottomrule
\end{tabular}
\end{table}

\textbf{Baseline Insights:}
\begin{itemize}
\item \primarytext{iCEpsilonGreedy Near-Oracle}: Achieves 93.2\% under optimal conditions, approaching theoretical maximum, outperforming neural baseline by 1.9\%.
\item \neural{iCPursuitNeuralUCB Excellence}: Achieves 91.1\% baseline with exceptional consistency (CV = 0.3\%), demonstrating strong performance in failure-free environments.
\item \primarytext{Exceptional Consistency}: iCEpsilonGreedy baseline CV of 0.4\% demonstrates remarkable stability, closely matching iCPursuitNeuralUCB stability (0.3\%).
\item \legacy{EXPNeuralUCB Baseline}: Achieves 86.3\% baseline, notably lower than top methods but with outstanding consistency (CV = 0.1\%).
\item \neural{GNeuralUCB Baseline}: Achieves 84.4\% baseline with solid consistency (CV = 0.5\%), demonstrating reliable operation in failure-free environments.
\end{itemize}

\subsection{Degradation Analysis: Stochastic vs. Baseline}

\begin{table}[H]
\centering
\caption{Oracle Efficiency Degradation Under Stochastic Failures}
\label{tab:degradation-analysis}
\begin{tabular}{@{}lcccc@{}}\toprule
\textbf{Algorithm} & \textbf{Baseline (\%)} & \textbf{Stochastic (\%)} & \textbf{Gap (\%)} & \textbf{Degradation (\%)} \\\midrule
\neural{iCPursuitNeuralUCB}   & \neural{91.1} & \neural{88.9} & \neural{2.2} & \neural{2.4\%} \\
\primarytext{iCEpsilonGreedy}   & \primarytext{93.2} & \primarytext{86.8} & \primarytext{6.4} & \primarytext{6.9\%} \\
\legacy{EXPNeuralUCB}           & \legacy{86.3} & \legacy{87.0} & \legacy{-0.7} & \legacy{-0.8\%} \\
\neural{GNeuralUCB}            & \neural{84.4} & \neural{85.2} & \neural{0.8} & \neural{0.9\%} \\
iCPursuit                       & 69.2 & 68.5 & 0.7 & 1.0\% \\
iCThompsonSampling              & 69.9 & 66.3 & 3.7 & 5.3\% \\
iCEXP4                          & 38.7 & 37.1 & 1.6 & 4.1\% \\
iCEpochGreedy                   & 39.3 & 37.6 & 1.7 & 4.3\% \\\bottomrule
\end{tabular}
\end{table}

\textbf{Degradation Interpretation:}
\begin{itemize}
\item \neural{iCPursuitNeuralUCB Resilience}: Shows minimal degradation (2.4\%) from baseline (91.1\%) to stochastic (88.9\%), maintaining superior absolute performance while recovering gracefully under failure conditions.
\item \primarytext{Relative Impact}: iCEpsilonGreedy's 6.9\% relative degradation is larger than secondary-tier algorithms due to its higher baseline, but maintains superior absolute performance to all pure methods.
\item \neural{GNeuralUCB Robustness}: Shows minimal degradation (0.9\%), demonstrating robust performance under failure conditions with stable baseline.
\item \success{Neural Advantage}: iCPursuitNeuralUCB's 2.2\% absolute gap significantly outperforms iCEpsilonGreedy's 6.4\% gap while maintaining a high baseline, representing a stronger robustness profile.
\end{itemize}

\subsection{Robustness Metrics (Stochastic Workload)}

\begin{table}[H]
\centering
\caption{iCMAB and Neural Robustness Characterization}
\label{tab:robustness}
\begin{tabular}{@{}lccccc@{}}\toprule
\textbf{Algorithm} & \textbf{Mean (\%)} & \textbf{Std Dev} & \textbf{Min (\%)} & \textbf{Max (\%)} & \textbf{Range} \\\midrule
\neural{iCPursuitNeuralUCB}   & \neural{88.9} & \neural{0.89} & \neural{87.8} & \neural{89.5} & \neural{1.7} \\
\success{iCEpsilonGreedy}   & \success{86.8} & \success{1.99} & \success{84.0} & \success{88.4} & \success{4.4} \\
\legacy{EXPNeuralUCB}       & \legacy{87.0} & \legacy{0.33} & \legacy{86.6} & \legacy{87.4} & \legacy{0.8} \\
\neural{GNeuralUCB}        & \neural{85.2} & \neural{0.46} & \neural{84.8} & \neural{85.4} & \neural{0.6} \\
iCPursuit                   & 68.5 & 4.46 & 62.3 & 72.6 & 10.3 \\
iCThompsonSampling          & 66.3 & 1.92 & 64.0 & 68.7 & 4.7 \\
iCEXP4                      & 37.1 & 0.29 & 36.7 & 37.4 & 0.7 \\
iCEpochGreedy               & 37.6 & 0.22 & 37.3 & 37.8 & 0.5 \\\bottomrule
\end{tabular}
\end{table}

\textbf{Consistency Analysis:}
\begin{itemize}
\item \neural{iCPursuitNeuralUCB Superior Stability}: Range of 1.7 percentage points with Std Dev 0.89 demonstrates exceptional consistency, superior to iCEpsilonGreedy (4.4 pp range) while maintaining higher absolute performance.
\item \success{iCEpsilonGreedy Stability}: Range of 4.4 percentage points with Std Dev 1.99 confirms tight, predictable performance among pure iCMAB algorithms.
\item \legacy{EXPNeuralUCB Stability}: Range of 0.8 pp with Std Dev 0.33 demonstrates exceptional consistency, best among neural alternatives.
\item \neural{GNeuralUCB Consistency}: Exceptional range of 0.6 pp with Std Dev 0.46 confirms outstanding run-to-run stability, best neural consistency.
\item \warning{iCPursuit Volatility}: Range of 10.3 percentage points (62.3\% to 72.6\%) indicates unreliable behavior despite pure-algorithm second-place ranking.
\end{itemize}

\section{Capacity-Scale Robustness}

\subsection{Performance Across Capacity Scales}

\begin{table}[H]
\centering
\caption{iCMAB and Neural Performance by Capacity Scale (1$\times$, 1.5$\times$, 2$\times$)}
\label{tab:capacity-scale}
\begin{tabular}{@{}lcccc@{}}\toprule
\textbf{Algorithm} & \textbf{1$\times$ (\%)} & \textbf{1.5$\times$ (\%)} & \textbf{2$\times$ (\%)} & \textbf{Max Variance (\%)} \\\midrule
\neural{iCPursuitNeuralUCB}   & \neural{88.6} & \neural{89.2} & \neural{88.9} & \neural{0.6} \\
\primarytext{iCEpsilonGreedy}   & \primarytext{86.8} & \primarytext{87.6} & \primarytext{86.1} & \primarytext{1.5} \\
\legacy{EXPNeuralUCB}           & \legacy{86.9} & \legacy{87.2} & \legacy{87.0} & \legacy{0.3} \\
\neural{GNeuralUCB}            & \neural{85.2} & \neural{85.3} & \neural{85.1} & \neural{0.2} \\
iCPursuit                       & 68.5 & 69.3 & 66.7 & 2.6 \\
iCThompsonSampling              & 66.3 & 67.0 & 66.1 & 0.9 \\
iCEXP4                          & 37.1 & 37.3 & 37.0 & 0.3 \\
iCEpochGreedy                   & 37.6 & 37.5 & 37.2 & 0.4 \\\bottomrule
\end{tabular}
\end{table}

\textbf{Capacity Scaling Insights:}
\begin{itemize}
\item \neural{iCPursuitNeuralUCB Robustness}: Superior 0.6\% max variance across capacity scales demonstrates exceptional configuration-independence and fundamental stability.
\item \primarytext{iCEpsilonGreedy Scaling}: 1.5\% max variance demonstrates dominance is not configuration-dependent among pure methods.
\item \neural{GNeuralUCB Superior Scaling}: 0.2\% max variance demonstrates best capacity robustness overall, superior configuration-independence.
\item \legacy{EXPNeuralUCB Stability}: 0.3\% max variance across scales confirms exceptional robustness independent of capacity regime.
\item \secondarytext{Consistent Ranking}: All capacity scales maintain consistent algorithm ranking, confirming fundamental performance hierarchy.
\end{itemize}

\section{Multi-Run Statistical Validation}

\subsection{3-run vs. 5-run Protocol Consistency}

\begin{table}[H]
\centering
\caption{Run-Count Stability (Stochastic Conditions)}
\label{tab:run-count}
\begin{tabular}{@{}lccc@{}}\toprule
\textbf{Algorithm} & \textbf{3-run Avg (\%)} & \textbf{5-run Avg (\%)} & \textbf{Difference (\%)} \\\midrule
\neural{iCPursuitNeuralUCB}   & \neural{88.9} & \neural{89.0} & \neural{0.1} \\
\success{iCEpsilonGreedy}   & \success{86.8} & \success{87.2} & \success{0.4} \\
\legacy{EXPNeuralUCB}       & \legacy{87.0} & \legacy{87.1} & \legacy{0.1} \\
\neural{GNeuralUCB}        & \neural{85.2} & \neural{85.3} & \neural{0.1} \\
iCPursuit                   & 68.5 & 68.2 & 0.3 \\
iCThompsonSampling          & 66.3 & 66.5 & 0.2 \\
iCEXP4                      & 37.1 & 37.0 & 0.1 \\
iCEpochGreedy               & 37.6 & 37.3 & 0.3 \\\bottomrule
\end{tabular}
\end{table}

\textbf{Statistical Conclusion:}
\begin{itemize}
\item \success{Negligible Differences}: All differences $\leq$0.4\%, confirming 3-run protocols capture representative behavior.
\item \neural{iCPursuitNeuralUCB Exceptional Convergence}: 0.1\% difference confirms outstanding stability across run-count variations.
\item \neural{GNeuralUCB Exceptional Convergence}: 0.1\% difference confirms neural consistency across run counts.
\item \success{Ranking Stability}: Both protocols maintain identical algorithm rankings.
\end{itemize}

\section{Implementation Gap Analysis}

\subsection{Current vs. Optimal Configuration}

\begin{table}[H]
\small
\centering
\caption{Implementation Gap Analysis: Current iCPursuit vs. Optimal Variants}
\label{tab:implementation-gap}
\begin{tabular}{@{}lcccc@{}}\toprule
\textbf{Metric} & \textbf{Current (iCPursuit)} & \textbf{Pure Optimal (iCEpsilonGreedy)} & \textbf{Neural Optimal} & \textbf{Best Gain} \\\midrule
\warning{Stochastic Efficiency} & 68.5\% & 86.8\% & \neural{88.9\%} & \neural{+20.4\%} \\
\warning{Stochastic Variance} & 6.5\% & 2.3\% & \neural{1.0\%} & \neural{6.5$\times$} \\
\warning{Baseline Efficiency} & 69.2\% & 93.2\% & \neural{91.1\%} & \neural{+21.9\%} \\
\warning{Min Performance} & 62.3\% & 84.0\% & \neural{87.8\%} & \neural{+25.5\%} \\
\warning{Max Performance} & 72.6\% & 88.4\% & \neural{89.5\%} & \neural{+16.9\%} \\\midrule
\warning{Degradation Gap} & 0.7\% & 6.4\% & \neural{2.2\%} & \neural{-5.7\%} \\\bottomrule
\end{tabular}
\end{table}

\textbf{Gap Summary:}
\begin{itemize}
\item \neural{Primary Opportunity (Neural Path)}: 20.4\% absolute performance gain by switching from iCPursuit to iCPursuitNeuralUCB baseline.
\item \primarytext{Primary Opportunity (Pure Path)}: 18.3\% absolute performance gain by switching from iCPursuit to iCEpsilonGreedy baseline.
\item \neural{Stability Improvement (Neural)}: 6.5$\times$ better stability (CV: 6.5\% $\rightarrow$ 1.0\%) makes iCPursuitNeuralUCB extraordinarily predictable for deployment.
\item \warning{Trade-off Resolution}: iCPursuitNeuralUCB achieves superior degradation resilience (2.2\% vs. 6.4\%), solving the robustness-performance trade-off that affects iCEpsilonGreedy.
\item \success{Net Benefit}: iCPursuitNeuralUCB provides 20\%+ performance advantage with superior stability, representing best overall value.
\end{itemize}

\section{Selection Framework for Hybrid Development}

\subsection{Performance Tiers}

Based on stochastic efficiency and cross-run consistency:

\begin{itemize}
\item \textbf{Tier 1 (Elite Neural)}: iCPursuitNeuralUCB (88.9\% stochastic, CV = 1.0\%) --- primary choice for maximum performance and stability.
\item \textbf{Tier 1 (Elite Pure / Neural Alternatives)}: iCEpsilonGreedy (86.8\% stochastic, CV = 2.3\%), EXPNeuralUCB (87.0\% stochastic, CV = 0.4\%) --- premier pure and alternative neural options.
\item \textbf{Tier 1-alt (Competitive Neural)}: GNeuralUCB (85.2\% stochastic, CV = 0.8\%) --- viable Gaussian-style neural alternative with exceptional consistency.
\item \textbf{Tier 2 (Secondary)}: iCPursuit (68.5\%, CV = 6.5\%), iCThompsonSampling (66.3\%, CV = 2.9\%) --- fallback options for ablations.
\item \textbf{Tier 3 (Lower-Bound)}: iCEXP4 (37.1\%), iCEpochGreedy (37.6\%) --- baselines only, insufficient for primary deployment.
\end{itemize}

\subsection{Hybrid Integration Criteria}

\begin{table}[H]
\footnotesize
\centering
\caption{Algorithm Selection Guidelines for Hybrid Architectures}
\begin{tabular}{@{}lll@{}}\toprule
\textbf{Use Case} & \textbf{Primary Choice} & \textbf{Rationale} \\\midrule
Primary Hybrid Baseline (Best) & iCPursuitNeuralUCB & 88.9\% efficiency, 1.0\% CV, superior stability \& degradation \\
Primary Hybrid Baseline (Pure) & iCEpsilonGreedy & 86.8\% efficiency, 2.3\% CV, best pure iCMAB \\
Neural Alternative 1 & EXPNeuralUCB & 87.0\% efficiency, 0.4\% CV, superior consistency \\
Neural Alternative 2 & GNeuralUCB & 85.2\% efficiency, 0.8\% CV, best capacity robustness \\
Hybrid Performance Target (Neural) & $\geq$88.9\% & Match or exceed iCPursuitNeuralUCB baseline \\
Hybrid Performance Target (Pure) & $\geq$86.8\% & Match or exceed iCEpsilonGreedy baseline \\
\bottomrule
\end{tabular}
\end{table}

\textbf{Concrete Deployment Recommendations:}
\begin{itemize}
\item \textbf{Primary Efficiency Threshold}: $\geq$88.9\% for neural-integrated hybrid acceptance (iCPursuitNeuralUCB parity).
\item \textbf{Secondary Efficiency Threshold}: $\geq$86.8\% for pure iCMAB hybrid acceptance (iCEpsilonGreedy parity).
\item \textbf{Stability Requirement}: CV $\leq$1.5\% for production deployment (matching iCPursuitNeuralUCB excellence).
\item \textbf{Degradation Profile}: Prefer algorithms with $\leq$2.5\% baseline-to-stochastic gap (superior to iCEpsilonGreedy's 6.4\%).
\item \textbf{Capacity Robustness}: GNeuralUCB at 0.2\% max variance provides best configuration-independence.
\item \textbf{Multi-Run Validation}: Require at least 3-run protocol for statistical confidence.
\item \textbf{Neural First Strategy}: iCPursuitNeuralUCB is recommended as primary baseline due to superior performance, stability, and degradation resilience.
\end{itemize}

\section{Conclusions}

\subsection{Key Findings}

\begin{enumerate}
\item \neural{iCPursuitNeuralUCB establishes new leadership} with 88.9\% stochastic efficiency, 91.1\% baseline performance, and strong consistency (CV = 1.0\%), demonstrating neural integration transforms iCPursuit performance.
\item \primarytext{iCEpsilonGreedy dominates pure iCMAB algorithms} with 86.8\% stochastic efficiency, 93.2\% baseline performance, and excellent consistency (CV = 2.3\%), confirming informed epsilon-greedy superiority.
\item \neural{Multiple neural pathways viable}: EXPNeuralUCB (87.0\%) and GNeuralUCB (85.2\%) establish competitive neural alternatives with distinct strengths.
\item \secondarytext{Clear tiered hierarchy} established: Tier 1 iCPursuitNeuralUCB (88.9\% neural), iCEpsilonGreedy (86.8\% pure), EXPNeuralUCB (87.0\% neural); Tier 1-alt GNeuralUCB (85.2\%); Tier 2 iCPursuit/iCThompsonSampling; Tier 3 iCEXP4/iCEpochGreedy.
\item \success{Neural Enhancement Dramatic}: iCPursuitNeuralUCB achieves 20.4\% improvement over pure iCPursuit (68.5\% $\rightarrow$ 88.9\%), demonstrating exceptional neural integration value for pursuit-based methods.
\item \neural{GNeuralUCB Consistency Excellence}: Achieves best capacity robustness (0.2\% max variance) and exceptional stochastic consistency (0.8\% CV), despite lower absolute performance (85.2\%).
\item \success{Robustness confirmed} across capacity scales (1$\times$--2$\times$), with iCPursuitNeuralUCB variance $<$0.6\%, demonstrating fundamental superiority independent of configuration.
\item \success{Statistical validation} shows 3-run and 5-run protocols differ by $\leq$0.4\%, confirming representativeness of current evaluation methodology.
\item \warning{Critical implementation gap identified}: Current hybrid uses iCPursuit (CMAB winner) instead of iCPursuitNeuralUCB (iCMAB+neural winner), losing 20.4\% performance.
\item \success{Immediate remediation available}: Implementing iCPursuitNeuralUCB requires 1--2 days engineering effort for significant performance gain.
\item \success{Degradation Resilience}: iCPursuitNeuralUCB's 2.2\% baseline-to-stochastic gap significantly outperforms iCEpsilonGreedy's 6.4\% gap, improving the robustness trade-off.
\end{enumerate}

\subsection{Recommendations for Hybrid Development}

\textbf{Immediate Actions (Priority Order):}
\begin{enumerate}
\item Implement iCPursuitNeuralUCB in primary hybrid framework position (best performance, stability, degradation).
\item Run convergence validation (3-run stochastic protocol) for iCPursuitNeuralUCB against current iCPursuit implementation.
\item Document the 20.4\% performance delta and neural integration strategy in the hybrid paper.
\item Prepare fallback evaluation of iCEpsilonGreedy + GNeuralUCB if neural integration encounters technical challenges.
\end{enumerate}

\textbf{Deployment Strategy:}
\begin{itemize}
\item \textbf{Primary Baseline}: iCPursuitNeuralUCB (88.9\%) for optimal hybrid performance, stability, and robustness.
\item \textbf{Fallback Baseline}: iCEpsilonGreedy (86.8\%) for pure iCMAB deployment if neural integration delayed.
\item \textbf{Performance Targets}: Hybrids must achieve $\geq$88.9\% stochastic efficiency to justify deployment.
\end{itemize}

\end{document}